% ============================================================
%  ch04b：§4.3 非共振、§4.4 定理 4.3.3.1 的证明（原文 pp.27 末–p.31 中段）
%  起止：自 “4.3. Non-resonance” 标题起，至 §4.4.6 末尾 “as was desired” 止；
%        其后 “4.5. Setup for the proof of Theorem 3.5.2.1” 起属 ch04c。
% ============================================================
\subsection{非共振}

\pnum{定义集合 $\mathcal{R}_X \subseteq \e^\times$ 为（几何）Frobenius 作用在 $H^1_{\et}(X,\e) \times H^2_{\et}(X,\e)$ 上的特征值所构成的集合。}

\begin{remark}
由曲线的 Weil 猜想，$1 \notin \mathcal{R}_X$。此外，$q$ 总是落在 $\mathcal{R}_X$ 之中（不过这一点与我们眼下的关切关系不大）。
\end{remark}

\pnum{}令 $H$ 为仿射代数群。

令 $V$ 为 $H$ 的有限维表示，并令 $\rho_V : H \to GL(V)$ 为相应的同态。令 $\ch_V : H^{ad}/H \times \mathbf{A}^1 \to \mathbf{A}^1$ 为使下图交换的映射：
\[
\begin{tikzcd}[column sep=large, row sep=large]
  H \times \mathbf{A}^1 \arrow[r, "\rho_V \times \id"] \arrow[d]
    & GL(V) \times \mathbf{A}^1 \arrow[d, "{(g,\lambda) \mapsto \det(g-\lambda\cdot\id_V)}"] \\
  H^{ad}/H \times \mathbf{A}^1 \arrow[r, "\ch_V"]
    & \mathbf{A}^1.
\end{tikzcd}
\]
显式地，对 $h \in H$ 与 $\lambda \in \mathbf{A}^1$，$\ch_V([h],\lambda)$ 就是 $\rho_V(h)$ 的特征多项式在 $\lambda$ 处的取值。

我们于是令 $(H^{ad}/H)^{\nonres} \subseteq H^{ad}/H$ 表示由满足 $\prod_{\lambda \in \mathcal{R}_X} \ch_{\mathfrak{h}}([h],\lambda) \neq 0$ 的共轭类 $[h]$ 构成的开集，其中 $\mathfrak{h}$ 是 $H$ 的伴随表示。显式地，$(H^{ad}/H)^{\nonres}$ 是那些使矩阵 $\rho_{\mathfrak{h}}(h) \in GL(\mathfrak{h})$ 在 $\mathcal{R}_X$ 中没有特征值的共轭类 $[h]$ 的集合。我们指出，开浸入 $(H^{ad}/H)^{\nonres} \hookrightarrow H^{ad}/H$ 显然是仿射的。

\begin{remark}
由注 4.3.1.1，$[1] \in (H^{ad}/H)^{\nonres}$。
\end{remark}

\begin{notation}
对任何配有到 $(H^{ad}/H)^{\nonres}$ 的结构映射的叠 $\mathcal{Y}$，我们令 $\mathcal{Y}^{\nonres} := \mathcal{Y} \times_{H^{ad}/H} (H^{ad}/H)^{\nonres}$。我们特别在 $\mathcal{Y} = \LS^{\arthm,\circ}_H$、配以结构映射 $\chi_H$ 的情形使用这个记号。
\end{notation}

\pnumt{主要几何结果。}{}\ 下面的结果把算术局部系统与伴随商作比较：

\begin{thm}
来自 (4.3) 的映射
\[
  \tau_H : (H^{ad}/H)^{\nonres} \to \LS^{\arthm,\circ,\nonres}_H
\]
是同构。
\end{thm}

该结果的证明是 §4.4 的内容。

\subsection{定理 4.3.3.1 的证明}

\pnumt{映射是同构的判别准则。}{}\ 我们首先作如下观察：

\begin{lem}
设 $f : \mathcal{Y}_1 \to \mathcal{Y}_2$ 是局部几乎有限型的代数叠（在代数闭域 $\e$ 上）之间的态射。

则 $f$ 是同构，当且仅当：
\begin{enumerate}[label=(\roman*)]
  \item $f$ 是形式 étale 的，即其余切复形 $\Omega^1_{\mathcal{Y}_1/\mathcal{Y}_2} \in \QCoh(\mathcal{Y}_1)$ 为零。
  \item 映射 $\mathcal{Y}_1(\e) \to \mathcal{Y}_2(\e)$ 是 (1-) 群胚的同构。
\end{enumerate}
这里我们要明确指出，条件 (ii) 可以拆成如下两个独立的条件：
\begin{enumerate}[label=(\text{ii}$_{\arabic*}$)]
  \item 对每个 $y_1 \in \mathcal{Y}_1(\e)$，自同构群的映射 $\Aut_{\mathcal{Y}_1(\e)}(y_1) \to \Aut_{\mathcal{Y}_2(\e)}(f(y_1))$ 是同构\footnote{我们强调，这里没有任何导出（derived）的余地：这是一对集合之间的映射。}。
  \item 对每个 $y_2 \in \mathcal{Y}_2(\e)$，存在 $y_1 \in \mathcal{Y}_1(\e)$ 以及同构 $f(y_1) \simeq y_2 \in \mathcal{Y}_2(\e)$。
\end{enumerate}
\end{lem}

\begin{proof}
只需证明：对每个局部几乎有限型、并配有映射 $S \to \mathcal{Y}_2$ 的仿射 $S$，映射 $S \times_{\mathcal{Y}_2} \mathcal{Y}_1 \to S$ 是同构。性质 (i) 与 (ii) 显然在这样的基变换下保持，于是我们可以假定 $\mathcal{Y}_2$ 是仿射概形。此外，众所周知 $\mathcal{Y}_1 \to \mathcal{Y}_2$ 是同构当且仅当 $\mathcal{Y}_1 \times_{\mathcal{Y}_2} \mathcal{Y}_2^{\mathrm{cl}} \to \mathcal{Y}_2^{\mathrm{cl}}$ 是同构；因此我们可以进一步假定 $\mathcal{Y}_2$ 是经典的。

此时 $\mathcal{Y}_1$ 是 $\e$-点上自同构群平凡的代数叠，从而是代数空间。又因 $\mathcal{Y}_1 \to \mathcal{Y}_2$ 是 étale 的，$\mathcal{Y}_1$ 也是经典的。现在 $f$ 是 radicial 态射（因为它局部有限型、且在 $\e$-点上单射），又是 étale 的，从而是开浸入。最后，由于 $f$ 在 $\e$-点上满射，它必定是同构。
\end{proof}

下面我们将对定理 4.3.3.1 中所考虑的映射 $\tau_H$ 验证上述这些假设。

\pnumt{约定、公式与符号。}{}\ 在继续之前，我们先确立某些重要的符号约定。粗略地说，要在概念上区分 $\tau_H([h])$ 与 $\tau_H([h^{-1}])$ 是困难的，但定理 4.3.3.1 确实能把它们区分开来，因此我们必须稍微更显式地解释该如何理解映射 $\tau_H$。（读者不妨跳过这段题外话，需要时再回头看。）

下面我们令 $\Frob_X : X \to X$ 表示几何 Frobenius 映射。

我们把 $X$ 上的 lisse Weil 层约定为如下的一对 $(\sigma,\alpha)$：其中 $\sigma \in \mathrm{Lisse}(X)^\heartsuit = \qLisse(X)^\heartsuit$，并配有同构 $\alpha : \sigma \xrightarrow{\ \sim\ } \Frob_X^*(\sigma)$。映射 $\alpha$ 的方向正是这里所说的“符号”。我们首先解释：为什么（以及在什么意义上）这个符号对于我们已有的约定而言是正确的。

注意 $V := x_0^*(\sigma)$ 是 $\pi_1^{\et}(X)$ 的一个表示 $\rho^{\mathrm{geom}}$。我们还得到一个同构
\begin{equation}
  V = x_0^*(\sigma) \xrightarrow{\ \alpha_{x_0}\ } \Frob_X(x_0)^*(\sigma) \xrightarrow{\ (4.2)\ } x_0^*(\sigma) = V,
\end{equation}
我们把它记作 $\rho(F)$。按 §4.1.1 的记号，可以同义反复地得到 $\rho(F) \circ \rho^{\mathrm{geom}}(g) = \rho^{\mathrm{geom}}(\pi_1^{\et}(\Frob_X)(g)) \circ \rho(F)$。由 (4.1)，我们可以把这个等式改写为
\[
  \rho(F)\rho^{\mathrm{geom}}(g)\rho(F)^{-1} = \rho^{\mathrm{geom}}(FgF^{-1}),
\]
于是我们得到 $\mathcal{W}_X$ 在 $V$ 上的表示，其中 $F$ 按 (4.4) 作用——倘若 $\alpha$ 的方向反了过来，我们就需要把 (4.4) 取逆。

类似地，对仿射代数群 $H$，一个 Weil $H$-局部系统指的是 $X$ 上的 $H$-局部系统 $\sigma_H$，连同同构 $\alpha : \sigma_H \xrightarrow{\ \sim\ } \Frob_X^*(\sigma_H)$（$H$-局部系统之间的）。作为推论，对 $h \in H$，$\tau_H([h])$ 的 $\sigma_H$ 是平凡的，而 $\alpha$ 由乘以 $h$ 给出。这就保证相应的 Weil 群表示 $\mathcal{W}_X \to H(\e)$ 通过 $\ZZ = \mathcal{W}_X/\pi_1^{\et}(X)$ 分解，并把生成元送到 $h$（正如它本来应该的那样）。

最后，对 lisse Weil 层 $(\sigma,\alpha)$，其上同调上自然的“几何” Frobenius 作用 $\phi_\sigma$ 由如下算子给出：
\begin{equation}
  C_{\et}(X,\sigma) \to C_{\et}(X,\Frob_X^*(\sigma)) \xrightarrow{\ \alpha^{-1}\ } C_{\et}(X,\sigma),
\end{equation}
其中第一个映射是同义反复的那个。

\begin{remark}
我们希望把上述材料的逻辑地位讲清楚。首先，我们论证了映射 $\alpha$ 应当按某个确定的方向来考虑。但这在某种程度上只是直觉性的论证，而非数学论证。毋宁说，我们只是把先前隐含的某个\footnote{说得更清楚些：$H^{ad}/H = (\mathbf{B}H)^{S^1}$ 有一个“回路反转”（loop reversal）的自同构，而我们需要消除这个自同构所带来的含糊。}约定明确化了（并说明它如何导向注 4.2.2.1 中非正式地暗示的那个定向）。从逻辑上讲，确立这一约定对于定理 4.3.3.1 的陈述是绝对必要的。
\end{remark}

\pnumt{$\tau_H$ 的形式 étale 性。}{}\ 我们将证明 $\tau_H|_{(H^{ad}/H)^{\nonres}}$ 是形式 étale 的。

首先注意，我们归结为验证切复形为零（例如，两侧的余切复形都是完美的）。此外，由于两侧都局部几乎有限型，我们可以在所有 $\e$-点的纤维上作验证。

一般地，对 $\sigma \in \LS^{\arthm}_H(\e)$，切复形可以如下计算\footnote{见 [2] 命题 2.2.2、§24.5.1。}：
\[
  T_{\LS^{\arthm}_H,\sigma} = C_{\et}(X,\mathfrak{h}_\sigma[1])^{\ZZ} = \mathrm{Ker}\bigl(\id - \phi_\sigma : C_{\et}(X,\mathfrak{h}_\sigma[1]) \to C_{\et}(X,\mathfrak{h}_\sigma[1])\bigr).
\]
这里 $\mathfrak{h}_\sigma$ 是由 $\sigma$ 诱导出的 $X$ 上的伴随 Weil 局部系统，$C_{\et}(X,\mathfrak{h}_\sigma[1])$ 是它的 étale 上同调复形（差一个平移），而我们取的是关于 Frobenius $\phi_\sigma$（来自 $\sigma$ 上的 Weil 结构）的作用的 $\ZZ$-不变量。

对 $h \in H(\e)$，切复形 $T_{H^{ad}/H,[h]} \in \Vect$ 是同伦核（即平移后的锥）：
\[
  T_{H^{ad}/H,[h]} = \mathrm{Ker}(\id - \Ad_h : \mathfrak{h}[1] \to \mathfrak{h}[1]),
\]
其中 $\Ad_h$ 是 $h$ 在 $\mathfrak{h}$ 上的伴随作用；更自然地说，把 $H^{ad}/H = \mathrm{Maps}(\mathbf{B}\ZZ,\mathbf{B}H)$ 写出来，该公式给出的是 $\ZZ$ 在 $\mathfrak{h}[1] = T_{\mathbf{B}H,\Spec(k)}$ 上作用的 $\ZZ$-不变量，其中生成元按 $\Ad_h$ 作用。（一个推论：我们看到，在上述对 $T_{H^{ad}/H,[h]}$ 的描述中本可以用 $\Ad_{h^{-1}}$ 代替 $\Ad_h$。）

另一方面，我们有
\[
  C_{\et}(X,\mathfrak{h}_{\tau_H([h])}[1]) = C_{\et}(X,\e) \otimes \mathfrak{h}[1],
\]
其 Frobenius 为
\[
  \phi_{\tau_H([h])} = \phi_X \otimes \Ad_{h^{-1}},
\]
其中逆号的出现是因为 (4.5) 中出现了 $\alpha^{-1}$。

Frobenius 等变的映射 $\e = H^0(X,\e) \to C_{\et}(X,\e)$（Frobenius 在来源上平凡地作用）诱导出交换图
\[
\begin{tikzcd}[column sep=3.4cm, row sep=large]
  \mathfrak{h}[1] = \e \otimes \mathfrak{h}[1] \arrow[r, "{\id - \Ad_{h^{-1}}}" description] \arrow[d]
    & \mathfrak{h}[1] = \e \otimes \mathfrak{h}[1] \arrow[d] \\
  C_{\et}(X,\e) \otimes \mathfrak{h}[1] \arrow[r, "{\id - \phi_X \otimes \Ad_{h^{-1}}}" description]
    & C_{\et}(X,\e) \otimes \mathfrak{h}[1].
\end{tikzcd}
\]
沿着各行过渡到（同伦）核，就得到 $\tau_H$ 在 $[h]$ 处的微分\footnote{要看出这一点，把 $H^{ad}/H$ 看作 $\Spec(k)$ 上算术局部系统的模叠，再把上面关于 $\LS^{\arthm}_H$ 的讨论相应地用上去即可。}。

因此我们看到：$\tau_H$ 在 $[h]$ 处是形式 étale 的，当且仅当
\[
  \id - \phi_X \otimes \Ad_{h^{-1}} : \tau^{\geq 1}C_{\et}(X,\e) \otimes \mathfrak{h} \to \tau^{\geq 1}C_{\et}(X,\e) \otimes \mathfrak{h}
\]
是同构；等价地，当且仅当上同调上的诱导映射
\begin{align*}
  \id - \phi_X \otimes \Ad_{h^{-1}} &: H^1_{\et}(X,\e) \otimes \mathfrak{h} \to H^1_{\et}(X,\e) \otimes \mathfrak{h},\\
  \id - \phi_X \otimes \Ad_{h^{-1}} &: H^2_{\et}(X,\e) \otimes \mathfrak{h} \to H^2_{\et}(X,\e) \otimes \mathfrak{h}
\end{align*}
都是同构。显然这恰好在 $1$ 不是 $\phi_X \otimes \Ad_{h^{-1}}$ 的特征值时发生，而这又恰好在 $1$ 不能写成 $\lambda \cdot \mu$（其中 $\lambda$ 是 $\phi_X$ 的特征值、$\mu$ 是 $\Ad_{h^{-1}} = \Ad_h^{-1}$ 的特征值）时发生，也就是说，当 $\Ad_h$ 没有特征值落在 $\mathcal{R}_X$ 中时发生。这正是 $[h]$ 落在 $(H^{ad}/H)^{\nonres}$ 中的定义条件，于是我们得到断言。

\pnumt{稳定化子。}{}\ 接下来我们验证引理 4.4.1.1 中的条件 (ii$_1$)。事实上，这一步是显然的，我们在此也绝不会用到非共振的微妙之处。我们在这里把论证明确写出来：

设 $\sigma \in \LS^{\arthm}_H(\e)$。按定义，$\sigma$ 提升为连续的 Weil 群表示 $\rho : \mathcal{W}_X \to H(\e)$，该提升在共轭意义下良定义。此时 $\Aut_{\LS^{\arthm}_H(\e)}(\sigma)$ 是 $\rho$ 的像在 $H(\e)$ 中的稳定化子。

类似地，$H^{ad}/H$ 中的一个 $\e$-点提升为某个 $h \in H(\e)$，而 $\Aut_{H^{ad}/H(\e)}([h])$ 是 $h$ 的稳定化子。

现在对 $h \in H$，$\tau_H([h])$ 就是 Weil 群表示 $\mathcal{W}_X \to \ZZ \xrightarrow{1 \mapsto h} H(\e)$，其稳定化子显然与 $h$ 的稳定化子一致。

\pnumt{提升同构类：设定。}{}\ 最后我们验证条件 (ii$_2$)。设 $\sigma \in \LS^{\arthm,\circ,\nonres}_H(\e)$。我们把 $\sigma$ 提升为连续表示 $\rho : \mathcal{W}_X \to H(\e)$。令 $\rho_0 : \pi_1^{\et}(X) \to H(\e)$ 表示 $\rho$ 在几何基本群上的限制；我们的任务是证明 $\rho_0$ 是平凡的。

这是一个具体的线性代数问题；我们在下面把细节写清楚。我们沿用 §4.1.1 的记号（特别地，$F \in \mathcal{W}_X$ 以及 $g \mapsto {}^F g$）。此外，我们引入更多记号：
\begin{itemize}
  \item 令 $H_\circ \subseteq H$ 表示 $\mathrm{Image}(\rho_0)$ 的 Zariski 闭包\zy{原文这里用小圆圈作下标（$H_\circ$），全文一贯如此（如 $H_\circ^{\mathrm{ab}}$），并非上标零；译文按原样排作下标小圆圈。[来源：A/原文 pp.30--31 版面]}。
  \item 令 $\theta : H_\circ \xrightarrow{\ \sim\ } H_\circ$ 表示 $\rho(F)$ 的伴随作用，即 $\theta(h) = \Ad_{\rho(F)}(h)$。
\end{itemize}

在上述记号下，注意我们有
\begin{equation}
  \rho_0({}^F g) = \theta(\rho_0(g)) \qquad \text{对一切 } g \in \pi_1^{\et}(X).
\end{equation}
现在我们对假设再作如下几点观察。

第一，$\rho_0$ 定义了 $\LS^{\restr,\circ}_H \subseteq \LS^{\restr}_H$ 中的一个点，这意味着由 [2] 命题 3.7.2，$\rho_0$ 通过 $H$ 的一个幂幺子群分解。等价地说，$H_\circ$ 是幂幺的。

第二，注意非共振条件意味着 $\Ad_{\rho(F)} = \mathrm{Lie}(\theta) : \mathfrak{h} \to \mathfrak{h}$ 在 $\mathcal{R}_X$ 中没有特征值。

\pnumt{提升同构类：证明。}{}\ 在上述记号下，我们的任务是证明 $H_\circ$ 是平凡的。由幂幺性，只需证明其交换化 $H_\circ^{\mathrm{ab}}$ 是平凡的。令 $V := \mathrm{Lie}(H_\circ^{\mathrm{ab}})$；由于 $H_\circ^{\mathrm{ab}}$ 是交换幂幺群，我们滥用记号，把它与它的 Lie 代数（即与之相伴的 $\e$-概形）等同起来。

由函子性，$\theta$ 诱导出 $V$ 的一个自同构，我们仍记作 $\theta$。设 $V \neq 0$；则对转置 $\theta^\vee : V^\vee \to V^\vee$，存在特征向量 $\mu \in V^\vee$；我们令 $\lambda \in \e^\times$ 表示其特征值。注意由非共振假设，$\lambda \notin \mathcal{R}_X$。

于是我们得到一个连续同态
\[
\begin{tikzcd}[column sep=large, row sep=large]
  \pi_1^{\et}(X) \arrow[r, "\rho_0"] \arrow[rrr, bend right=22, "\bar{\rho}_0"']
    & H_\circ(\e) \arrow[r]
    & V \arrow[r, "\mu"]
    & \e
\end{tikzcd}
\]
由 (4.6)，它满足
\begin{equation}
  \bar{\rho}_0({}^F g) = \lambda \cdot \bar{\rho}_0(g).
\end{equation}
我们还指出，$\mathrm{Ker}\bigl(H_\circ(\e) \to V \to \e\bigr)$ 是 $H_\circ$ 的某个代数子群的 $\e$-点集，于是按后者的定义，同态 $\bar{\rho}_0$ 必须是非平凡的。

现在 $\bar{\rho}_0$ 延拓为 $\e$-线性映射 $H_{\et,1}(X,\e) \to \e$，即它来自一个非零的上同调类 $\eta \in H^1_{\et}(X,\e)$。由于 $g \mapsto {}^F g$ 在 $H_{\et,1}(X,\e)$ 上诱导出（几何）Frobenius（见 (4.1)），(4.7) 意味着
\[
  \phi_X(\eta) = \lambda \cdot \eta.
\]
这与非共振假设矛盾，于是我们断定 $V = 0$，此即所欲证者。
